
\documentclass[11pt,a4paper]{article}

% -------------------------------------------------
% Encoding and typography
% -------------------------------------------------
\usepackage[T1]{fontenc}
\usepackage[utf8]{inputenc}
\usepackage{microtype}
\usepackage{amsmath}
\usepackage{amssymb}

% -------------------------------------------------
% Page layout and tables
% -------------------------------------------------
\usepackage[
    top=2.5cm,
    bottom=2.5cm,
    left=2.7cm,
    right=2.7cm
]{geometry}
\usepackage{array}
\usepackage{ragged2e}
\usepackage{booktabs}
\usepackage[table]{xcolor}

% -------------------------------------------------
% Links
% -------------------------------------------------
\usepackage[hidelinks]{hyperref}

% -------------------------------------------------
% Paragraph formatting
% -------------------------------------------------
\setlength{\parindent}{0pt}
\setlength{\parskip}{0.6em}

% -------------------------------------------------
% Working-note helpers
% -------------------------------------------------
\definecolor{verifiedgreen}{RGB}{46,125,82}
\definecolor{reflectionorange}{RGB}{175,105,20}
\definecolor{questionblue}{RGB}{48,86,145}

\newcommand{\verified}{\textcolor{verifiedgreen}{\textbf{Verified observation.}}}
\newcommand{\reflection}{\textcolor{reflectionorange}{\textbf{Research reflection.}}}
\newcommand{\openquestion}{\textcolor{questionblue}{\textbf{Open question.}}}

\begin{document}

% =================================================
% Header
% =================================================
\begin{flushleft}

    {\LARGE\textbf{Geospatial Foundation Models Benchmarking}}\\[0.3em]

    {\large\color{black!65}
    Research Reflections and Open Methodological Questions
    }\\[1.1em]

    \textbf{Stefano Infusini}
    \hfill
    University of Genoa -- DIBRIS\\[-0.1em]

    {\small\color{black!65}
    MSc Computer Science -- Artificial Intelligence
    }
    \hfill
    {\small September 2026}

\end{flushleft}

\vspace{-0.6em}
\rule{\textwidth}{0.4pt}
\vspace{1.1em}

\section*{Purpose and status}

This is a \emph{living research-notes document}. Its purpose is to preserve observations,
criticisms, hypotheses and open methodological questions that emerge while studying
Geospatial Foundation Model (GFM) evaluation. It is deliberately not written as a final thesis
methodology section. Statements should progressively be verified against primary papers,
repositories and dataset documentation before being promoted to thesis claims.

A central working assumption is that a downstream benchmark score is not generated by
architecture alone. A more realistic abstraction is

\[
S =
f(
\text{architecture},
\text{pretraining objective},
\text{pretraining data},
\text{sensor},
\text{bands},
\text{resolution},
\text{geography},
\text{temporality},
\text{input adaptation},
\text{decoder},
\text{training regime}
).
\]

Consequently, a cross-model leaderboard can be useful while still having limited causal
interpretability. If two GFMs differ simultaneously in architecture, pretraining corpus,
spectral support, resolution and modality, a downstream ranking cannot by itself establish
which of those factors caused the performance difference.

% =================================================
\section{Metrics: IoU, mIoU and the foreground-class problem}

For a semantic class \(c\), Intersection over Union is

\[
\mathrm{IoU}_c
=
\frac{TP_c}{TP_c + FP_c + FN_c}.
\]

For a \(C\)-class segmentation problem, mean IoU is

\[
\mathrm{mIoU}
=
\frac{1}{C}\sum_{c=1}^{C}\mathrm{IoU}_c.
\]

For binary segmentation with a foreground class \(f\) and background class \(b\),

\[
\mathrm{mIoU}
=
\frac{\mathrm{IoU}_{b}+\mathrm{IoU}_{f}}{2}.
\]

\verified\ mIoU is a macro-average: each class receives equal weight regardless of its pixel
frequency. It is therefore substantially less sensitive to class imbalance than ordinary
pixel accuracy.

\reflection\ Even so, for highly asymmetric binary applications, mIoU may hide the quantity
that is operationally most interesting. If background IoU is both high and relatively stable,
variation in foreground IoU is compressed by the averaging operation. If
\(\mathrm{IoU}_{b}\) stays fixed, then

\[
\Delta\mathrm{mIoU}
=
\frac{1}{2}\Delta\mathrm{IoU}_{f}.
\]

For example, the purely hypothetical pair

\[
\mathrm{IoU}_{b}=0.98,\qquad \mathrm{IoU}_{f}=0.30
\]

gives

\[
\mathrm{mIoU}=0.64,
\]

whereas

\[
\mathrm{IoU}_{b}=0.98,\qquad \mathrm{IoU}_{f}=0.72
\]

gives

\[
\mathrm{mIoU}=0.85.
\]

The foreground performance has changed dramatically, but the aggregate metric partially
compresses that change.

\textbf{Important correction.} A reported mIoU value alone does \emph{not} allow the
foreground IoU to be reconstructed. For example, a mIoU of 83.62\% does not imply a specific
burn-scar IoU unless the background IoU (or the per-class confusion information) is also
known.

\reflection\ For binary EO segmentation tasks such as burn-scar or flood mapping, a more
transparent thesis evaluation would report at least

\[
\boxed{\mathrm{mIoU},\quad
\mathrm{IoU}_{\mathrm{foreground}},\quad
\mathrm{IoU}_{\mathrm{background}}}
\]

and, where useful, foreground precision, recall and F1/Dice. This retains benchmark
comparability while exposing whether a model is genuinely improving the target phenomenon.

% =================================================
\section{PANGAEA case study: HLS Burn Scars}

\verified\ PANGAEA evaluates GFMs over eleven downstream datasets and reports mIoU for
semantic segmentation and change detection, and RMSE for regression \cite{pangaea}. The HLS
Burn Scars dataset contains 804 HLS scenes of size \(512\times512\), with six bands
(Blue, Green, Red, NIR, SWIR1 and SWIR2). PANGAEA reports an approximate class distribution
of 11\% burn scar, 88\% not burned and 1\% no-data \cite{pangaea}.

Using 100\% of the downstream training data, PANGAEA Table 5 reports:

\begin{table}[h]
\centering
\small
\begin{tabular}{lr}
\toprule
Model & HLS Burn Scars mIoU (\%) \\
\midrule
UNet baseline & 84.51 \\
Prithvi & 83.62 \\
CROMA & 82.42 \\
S12-MAE & 81.91 \\
S12-Data2Vec & 81.91 \\
S12-DINO & 81.72 \\
S12-MoCo & 81.58 \\
GFM-Swin & 76.90 \\
Scale-MAE & 76.68 \\
RemoteCLIP & 76.59 \\
\bottomrule
\end{tabular}
\caption{Selected HLS Burn Scars results from PANGAEA Table 5.}
\end{table}

\verified\ The fully supervised UNet trained from scratch reaches 84.51 mIoU, higher than all
GFMs in the main 100\%-label comparison. In Table 5, UNet is among the top two models on six
of the eleven benchmark datasets, compared with three for Scale-MAE and two for CROMA and
DOFA \cite{pangaea}.

\verified\ PANGAEA explicitly argues that GFMs do not consistently outperform supervised
models. The authors observe that supervised baselines can be particularly competitive on
lower-resolution tasks and that models pretrained with broader spectral information tend to
perform better on tasks where spectral detail matters. They specifically note that Prithvi,
which was pretrained on HLS imagery, performs strongly on HLS Burn Scars \cite{pangaea}.

\reflection\ This is important for interpretation. The HLS Burn Scars result should not be
read simply as ``UNet is better than foundation models.'' It shows that, when full downstream
supervision is available, a task-specific convolutional model can be extremely competitive
on a relatively constrained dense-prediction problem. The result is consistent with the more
general possibility that pretraining is most valuable when downstream labels are scarce,
the task is difficult, or the representation must transfer across a meaningful domain gap.

\reflection\ A physically plausible explanation for burn-scar mapping is that the task is
strongly driven by spectral changes involving NIR and SWIR information. However, statements
about particular indices such as NBR/dNBR, or about Transformer patchification being a cause
of underperformance, should be treated as hypotheses unless directly supported by an
experiment or source. They should not be presented as conclusions of PANGAEA.

\openquestion\ PANGAEA reports aggregate mIoU in the main leaderboard. For the thesis, should
binary segmentation datasets additionally be re-evaluated using foreground-only IoU and
foreground F1/recall? This could reveal differences that are obscured by the aggregate score.

% =================================================
\section{Data scarcity changes the interpretation}

\verified\ PANGAEA also evaluates 10\%, 50\% and 100\% label regimes \cite{pangaea}. The
results are heterogeneous. On HLS Burn Scars, for example, SpectralGPT reports 83.35 mIoU at
10\% labels, while the UNet baseline reports 79.46. Other GFMs do not necessarily outperform
UNet in the same setting.

\reflection\ Therefore, the relevant scientific question is not simply whether a GFM wins
with full supervision. A more informative question is

\[
\boxed{
\text{How does performance degrade as the amount of geographically independent supervision decreases?}
}
\]

This is directly relevant to a thesis centred on label efficiency. A GFM may fail to win at
100\% labels but still provide a meaningful advantage in the low-label regime. Conversely, a
model that remains weak even when labels are scarce provides evidence that pretraining does
not automatically imply useful transfer.

% =================================================
\section{Benchmark fairness: heterogeneous models do not receive identical information}

\verified\ PANGAEA explicitly acknowledges that the selected GFMs expect different spectral
bands and sensors. Its band-adaptation procedure matches downstream bands to those expected
by each pretrained encoder; RGB-pretrained GFMs therefore use only RGB channels from
multispectral datasets, while unavailable expected bands can be zero-padded \cite{pangaea}.

The supervised UNet and ViT baselines, by contrast, are trained from scratch using all
available downstream bands \cite{pangaea}.

\reflection\ This protocol is defensible if the goal is to compare \emph{complete pretrained
systems under their native input constraints}. However, the models do not necessarily receive
the same information. A score difference can therefore reflect both representation quality
and differences in spectral access.

For example,

\[
\text{RGB-pretrained model}
\rightarrow
\{R,G,B\}
\]

may be compared on the same nominal benchmark with

\[
\text{multispectral model}
\rightarrow
\{R,G,B,NIR,SWIR,\ldots\}.
\]

Thus ``same dataset'' does not necessarily mean ``same effective input.''

\reflection\ A future thesis comparison should explicitly record an \emph{information-access
profile} for every model--dataset pair rather than hiding this mismatch inside preprocessing.

% =================================================
\section{Pretraining--downstream alignment versus pretraining exposure}

A key distinction is needed between two different phenomena.

\subsection*{1. Pretraining--downstream alignment}

This means that a model was pretrained on data that are similar to the downstream benchmark
in sensor, spectral bands, resolution, geography or acquisition characteristics.

For example,

\[
\text{Prithvi pretraining: HLS}
\qquad\Longrightarrow\qquad
\text{downstream: HLS Burn Scars}.
\]

\verified\ PANGAEA itself notes that such a match can be associated with improved downstream
performance \cite{pangaea}.

This is not automatically leakage. It is a legitimate property of the pretrained model.
However, it is a confounder if one attempts to interpret the leaderboard as a clean comparison
of architectures or self-supervised objectives.

\subsection*{2. Actual pretraining exposure}

A stronger concern is whether the pretraining corpus contains the same images, locations,
time periods, or near-duplicates used in a downstream benchmark:

\[
D_{\mathrm{pretrain}}\cap D_{\mathrm{benchmark}}\neq\varnothing.
\]

Even if downstream labels were never seen during self-supervised pretraining, exposure to the
same geography or acquisitions may make an apparent ``geographic generalization'' result
easier to interpret incorrectly.

\openquestion\ For every selected GFM and benchmark, can the pretraining corpus be audited
well enough to identify:
\begin{itemize}
    \item exact dataset overlap;
    \item geographic overlap;
    \item temporal overlap;
    \item sensor and spectral overlap;
    \item resolution similarity?
\end{itemize}

This should be treated as an empirical audit question rather than assuming that overlap exists.

% =================================================
\section{Proposed model--dataset compatibility taxonomy}

A useful working taxonomy for every model--dataset pair is:

\begin{center}
\begin{tabular}{p{3.7cm}p{9.0cm}}
\toprule
\textbf{Category} & \textbf{Meaning} \\
\midrule
Native-compatible &
Dataset supplies the modality and channels expected by the pretrained model with only ordinary normalization/resizing. \\

Deterministically adaptable &
A transparent, fixed transformation is needed, such as selecting matching bands or resizing. \\

Substantial adaptation required &
Missing modalities/bands, zero-padding, learned input projections, or major temporal/resolution transformations materially alter the input interface. \\

Not scientifically meaningful &
The required transformation is so severe that the result would no longer be a credible test of the released pretrained representation. \\
\bottomrule
\end{tabular}
\end{center}

\reflection\ Benchmark results should ideally be stratified by these categories. Otherwise,
a model that is evaluated natively and a model that is heavily adapted may appear as directly
comparable even though the experimental conditions are qualitatively different.

% =================================================
\section{So2Sat LCZ42: lessons from the classification case}

\verified\ So2Sat LCZ42 is a 17-class patch-classification benchmark for Local Climate Zones,
with 10 built classes and seven natural classes. The original benchmark contains 400,673
paired Sentinel-1/Sentinel-2 samples. Each sample covers \(32\times32\) pixels at 10\,m,
corresponding to approximately \(320\times320\) m \cite{so2sat}.

The labels originate from manually delineated LCZ polygons. A Landsat-8 Random Forest is used
during annotation to generate a dense citywide LCZ map that helps annotators identify
potential inconsistencies and missing training examples. The Random Forest is therefore best
understood as an annotation-debugging or consistency tool, not as an independent source of
ground truth. Final label quality is additionally assessed through human review and a
quantitative relabeling experiment with ten remote-sensing experts \cite{so2sat}.

\reflection\ This labeling pipeline illustrates an important benchmark-design principle:
model-assisted annotation can efficiently expose suspicious regions, but agreement between
the assisting model and the original annotator is not independent evidence of correctness,
because the model itself was trained from those annotations.

\verified\ The original machine-learning split trains on 32 main cities plus ten add-on
areas, while ten other cities are reserved for validation and testing. Within those held-out
cities, validation and test labels are geographically separated into different city halves
\cite{so2sat}.

\reflection\ This creates a useful cross-city evaluation, but validation and test still come
from the same city domains. Therefore, hyperparameter selection on validation data may reveal
information about the same city-level domains later used for testing. A stricter geographic
protocol could instead use disjoint training, validation and test cities.

% =================================================
\section{So2Sat and multimodal GFM compatibility}

\verified\ So2Sat supplies an unusual Sentinel-1 representation with eight derived real-valued
channels and ten Sentinel-2 bands \cite{so2sat}. The released CROMA implementation, in
contrast, expects exactly two Sentinel-1 backscatter channels and twelve Sentinel-2 channels
\cite{croma}.

Therefore,

\[
\text{So2Sat S1: }8\text{ channels}
\qquad\neq\qquad
\text{CROMA S1: }2\text{ channels},
\]

and

\[
\text{So2Sat S2: }10\text{ channels}
\qquad\neq\qquad
\text{CROMA S2: }12\text{ channels}.
\]

\reflection\ So2Sat can still be an excellent classification benchmark while being
\emph{non-native} for a particular GFM. Dataset quality and model compatibility are different
questions.

A possible deterministic adaptation for CROMA would be to select the So2Sat SAR channels
that correspond most closely to VV/VH intensity/backscatter and explicitly handle the missing
optical channels. However, this should be labelled as an adapted setting, not presented as
native CROMA input.

\openquestion\ Should a thesis benchmark force every model onto every dataset using generic
adaptation, or should it evaluate only scientifically defensible model--dataset pairs? The
second approach may produce a smaller but more interpretable benchmark.

% =================================================
\section{A more careful interpretation of GFM leaderboards}

The working interpretation should be:

\[
\boxed{
\text{Model A outperforming Model B on Dataset X}
}
\]

means

\[
\boxed{
\text{the complete pretrained system A transfers better under that experimental configuration}
}
\]

and does \emph{not} automatically mean

\[
\boxed{
\text{architecture A is intrinsically better than architecture B}.
}
\]

Differences in pretraining modality, spectral range, resolution, geography, data volume,
objective and input adaptation all remain potential explanations.

\reflection\ This suggests that a useful contribution of the thesis may be not merely another
leaderboard, but an evaluation protocol that makes these hidden compatibility and alignment
variables explicit.

% =================================================
\section{Benchmark difficulty: interpreting So2Sat classification performance}

\reflection\ An overall accuracy of 50--60\% on So2Sat LCZ42 should not be interpreted in
the same way as 50\% accuracy in a balanced binary classification problem. In a balanced
binary task, random guessing gives

\[
P(\mathrm{correct})=\frac{1}{2}=50\%.
\]

With 17 equally likely classes, uniform random guessing would instead give

\[
P(\mathrm{correct})=\frac{1}{17}\approx5.9\%.
\]

So a 50--60\% overall accuracy on So2Sat is far above a uniform-random baseline. Since the
real class distribution is not perfectly balanced, a more informative reference would also
include a majority-class baseline rather than relying only on \(1/17\).

\reflection\ So2Sat is intrinsically difficult for several reasons. First, some LCZ classes
are semantically close, for example compact versus open urban forms and low-rise versus
mid-rise categories. The human-validation experiment itself shows non-negligible confusion
between some of these classes. Second, the benchmark contains meaningful geographic
structure: the original evaluation includes cities that are held out from training. Third,
the LCZ target captures aspects of urban morphology rather than only simple surface-cover
categories. Therefore, an overall accuracy around 50--60\% should not automatically be
interpreted as evidence that the benchmark is poorly constructed.

\reflection\ For benchmarking purposes, this difficulty can be useful. A benchmark that is
nearly saturated at 98--99\% accuracy has limited ability to discriminate between models or
to reveal differences across label-efficiency and geographic-shift regimes. In contrast, a
non-saturated benchmark may preserve enough headroom to study

\[
100\%\text{ labels}\rightarrow10\%\rightarrow1\%
\]

and

\[
\mathrm{IID}\rightarrow
\mathrm{spatial\ holdout}\rightarrow
\mathrm{geographic\ OOD}.
\]

However, the opposite failure mode is also possible. If all models collapse to similarly low
performance, the benchmark may enter a practical floor regime in which differences between
models and label fractions become difficult to interpret.

\openquestion\ Before So2Sat is fixed as a final classification benchmark, modern GeoFMs
should be evaluated under a clearly specified protocol to determine whether the dataset lies
in a useful intermediate difficulty regime: neither saturated nor dominated by floor effects.

% =================================================
\section{GeoFM performance is task- and domain-dependent}

\reflection\ A recurring lesson from current GeoFM benchmarks is that the term
``foundation model'' should not be interpreted as implying uniformly superior performance
across Earth-observation downstream tasks. A model may transfer very well to one benchmark
and poorly to another.

A more realistic interpretation of observed downstream performance is

\[
S =
f(
\text{pretraining},
\text{sensor},
\text{bands},
\text{resolution},
\text{task},
\text{geography},
\text{decoder},
\text{adaptation regime}
).
\]

For example, a model pretrained on Sentinel-1/Sentinel-2 data may be relatively well aligned
with a downstream Sentinel benchmark, whereas a model pretrained primarily on RGB
high-resolution imagery may face a much larger input-domain shift when evaluated on
10\,m multispectral Sentinel-2 data.

Therefore, a result such as

\[
FM_A=80\%,\qquad FM_B=55\%
\]

does not by itself imply that \(FM_A\) has universally better representations. A more careful
statement is:

\begin{quote}
Under this sensor/task/domain configuration, \(FM_A\) transfers better than \(FM_B\); the
result should be interpreted in light of differences in pretraining modality, spectral
coverage, spatial resolution, geographic exposure and input adaptation.
\end{quote}

\reflection\ This distinction should be part of the thesis framework. The objective should
not be to produce only a single global ranking of GFMs, but to explain \emph{under which
conditions} different pretrained systems transfer well or poorly.

% =================================================
\section{Toward a model--task compatibility matrix}

A useful way to organize the experimental results is to combine downstream performance with
metadata about the compatibility of each model--dataset pair.

\[
\begin{array}{c|cccc}
 & \text{Classification}
 & \text{Segmentation}
 & \text{Regression}
 & \text{Change detection}\\
\hline
\text{CROMA} & ? & ? & ? & ?\\
\text{DOFA} & ? & ? & ? & ?\\
\text{Prithvi} & ? & ? & ? & ?\\
\text{Clay} & ? & ? & ? & ?
\end{array}
\]

For each cell, the benchmark record should additionally track:

\begin{itemize}
    \item whether the downstream input is native to the pretrained model;
    \item whether the downstream sensor matches the pretraining sensor;
    \item whether the spatial resolution / GSD is similar;
    \item whether the spectral bands overlap substantially;
    \item whether the downstream geography may have been seen during pretraining;
    \item whether deterministic or substantial adaptation is required.
\end{itemize}

\reflection\ This contextual matrix is potentially more informative than a plain leaderboard.
It makes explicit that model performance is conditional on the model--dataset relationship
rather than being an intrinsic scalar property of the model.

For So2Sat, the current working characterization is:

\[
\boxed{\text{Task: patch classification}}
\]

\[
\boxed{\text{Difficulty: substantial and non-saturated}}
\]

\[
\boxed{\text{Historical 50--60\% OA is far above uniform random guessing}}
\]

\[
\boxed{\text{Strong geographic structure}}
\]

but also

\[
\boxed{\text{FM compatibility and pretraining alignment require careful control}.}
\]

\reflection\ This last point may be one of the main ways to move beyond simply reproducing
a PANGAEA-style leaderboard: the benchmark should explicitly model compatibility,
pretraining alignment and possible exposure when interpreting performance.



% =================================================
\section{Candidate thesis-level methodological questions}

The following questions should remain open while the benchmark suite is being designed:

\begin{enumerate}
    \item Does apparent GFM sample efficiency survive when spatial leakage is removed and
    supervision is geographically independent?

    \item How much of cross-model performance can be explained by pretraining--downstream
    sensor, spectral and resolution alignment?

    \item Are conclusions about model superiority stable when analysis is restricted to
    native-compatible model--dataset pairs?

    \item For binary segmentation tasks, do foreground-only metrics change the ranking or
    interpretation obtained from aggregate mIoU?

    \item Does geographic OOD performance change when validation and test domains are also
    geographically disjoint, rather than merely spatially separated within the same domains?

    \item Can pretraining exposure be audited sufficiently to distinguish true geographic
    generalization from transfer to already-seen locations or acquisition regimes?
\end{enumerate}

% =================================================
\section{Working checklist for future benchmark analysis}

For each dataset:
\begin{itemize}
    \item What is the statistical unit: patch, pixel, object, field, event, city?
    \item How were labels created and independently validated?
    \item What is the official split and what kind of shift does it actually measure?
    \item What leakage mechanisms are possible?
    \item What metrics are reported, and do they hide important per-class behaviour?
    \item Which sensor, spectral bands, GSD and temporal information are available?
\end{itemize}

For each model--dataset pair:
\begin{itemize}
    \item Is the input native or adapted?
    \item Which downstream bands can the model actually use?
    \item Are missing expected bands zero-padded?
    \item How closely do pretraining sensor, spectrum and resolution match the benchmark?
    \item Is pretraining geographic or temporal exposure known?
    \item Is the encoder frozen, partially tuned or fully fine-tuned?
    \item Is the downstream head/decoder comparable across models?
\end{itemize}

% =================================================
\begin{thebibliography}{9}

\bibitem{pangaea}
V. Marsocci et al.,
\textit{PANGAEA: A Global and Inclusive Benchmark for Geospatial Foundation Models},
arXiv:2412.04204, 2024.
\href{https://arxiv.org/abs/2412.04204}{arXiv link}.

\bibitem{so2sat}
X. X. Zhu et al.,
\textit{So2Sat LCZ42: A Benchmark Data Set for the Classification of Global Local Climate Zones},
IEEE Geoscience and Remote Sensing Magazine, 8(3), 76--89, 2020.
DOI: 10.1109/MGRS.2020.2964708.

\bibitem{croma}
A. Fuller, K. Millard, and J. R. Green,
\textit{CROMA: Remote Sensing Representations with Contrastive Radar-Optical Masked Autoencoders},
NeurIPS, 2023.
\href{https://github.com/antofuller/CROMA}{Official repository}.

\end{thebibliography}

\end{document}
